\documentclass[12pt]{article}
\usepackage{amsmath, amssymb, enumitem, fancyhdr, graphicx, libertine, nth,
  pgfplots, tikz, xcolor}
\pgfplotsset{compat=1.11}
\usepackage[margin=1in]{geometry}
\usepackage[libertine]{newtxmath}
\pagestyle{fancy}
\definecolor{solution-color}{HTML}{000080}
\chead{\emph{EC201 -- Practice Midterm}}

% Define new topic command
\newcommand{\topic}[1]{\medskip\begin{center}\large\textsc{#1}\normalsize
  \end{center}}

% Define question counter:
\newcounter{question}[section]
\newcommand{\question}[1]{\refstepcounter{question}
  \subsubsection*{Question~\thequestion~-- #1}}

% Subquestions are enumerate with letters
\newenvironment{subquestions}{
  \begin{enumerate}[label=(\roman*)]}{\end{enumerate}}
\begin{document}
\question{Preferences (5 Points)}
The following are common assumptions about preferences:
\begin{itemize}
  \item Completeness
  \item Transitivity
  \item Monotonicity
  \item Convexity
\end{itemize}
Consider the following scenario:
\begin{itemize}
  \item I ask you which of the bundles $\left( x_1,x_2 \right)=\left( 2,8
    \right)$ and $\left(x_1,x_2  \right)=\left( 8, 2 \right)$ you prefer you
    tell me you are indifferent between the two.
  \item I ask you which of the bundles $\left( x_1,x_2 \right)=\left( 2,8
    \right)$ and $\left( x_1,x_2 \right)=\left( 5,5 \right)$ you prefer and you
    tell me that you strictly prefer $\left( x_1,x_2 \right)=\left( 2,8
    \right)$.
\end{itemize}
Which one of these assumptions do your choices violate and why?


\question{Choice and Demand (15 Points)}
Suppose your utility function for goods 1 and 2 was:
\[
  u\left( x_1,x_2 \right)=\max\left\{ x_1,x_2 \right\}
\]
\begin{subquestions}
  \item \textbf{[5 Points]} Sketch some indifference curves for this utility function.
  \item \textbf{[5 Points]} If $x_1=1$ and $x_2=2$, what is your $MRS$? Interpret the number.
  \item \textbf{[5 Points]} Find the demand functions $x_1\left( p_1,p_2,m \right)$ and $x_2\left(
    p_1,p_2, m\right)$
  % \item Suppose initially $p_1>p_2$. Then $p_2$ increases so much that we end
  %   up with $p_2>p_1$. Does the demand for good 1 increase or decrease? Are the
  %   goods substitutes or complements?
  % \item Can you come up with an example of two goods where you might have such
  %   preferences? Feel free to make assumptions about a particular scenario.
\end{subquestions}
\question{Income and Substitution Effects (10 Points)}
You live for one period. There are two goods with prices $p_1$ and $p_2$. You have income $m$ to spend on the two goods. Both goods are normal and your preferences satisfy all the main assumptions.

For each of the following parts, draw \textsl{\textbf{separate}} graphs.
\begin{subquestions}
\item \textbf{[2 Points]} Draw the budget line, your optimal consumption bundle and the indifference curve associated with the optimal choice. Label your axes and the points where the budget line intercepts the axes.
\item \textbf{[3 Points]} The price of good 1 falls. Describe what happens in a diagram. Show your new budget line, optimal bundle and indifference curve associated with the optimal choice.
\item \textbf{[5 Points]} Decompose the change in demand for good 1 into income and substitution effects in a diagram. Show the optimal choice you would make if you only experienced a relative change in prices.
\end{subquestions}

\question{Intertemporal Choice (5 Points)}
What is the most you would pay for a security that paid you \$100 in one year
from now and \$200 in two years from now? The interest rate is 5\%.
\question{Uncertainty (15 Points)}
A Roulette wheel has 38 numbered pockets, labelled 0, 00, 1, 2, 3, \dots, 34, 35,
36. There are 18 red numbers, 18 black numbers and 2 green numbers (0 and 00).
The wheel is spun and a ball is thrown. The ball is equally likely to land on
any of the 38 pockets.

You have \$100. If you bet \$100 on red and it lands on red, you get \$200
back (your profit is \$100). If it lands on black or green, you get nothing
(you lose \$100).
\begin{subquestions}
  \item \textbf{[5 Points]} What is the expected value of wealth from taking the bet?
  \item \textbf{[5 Points]} Suppose your utility function for wealth is $u\left( W
    \right)=\sqrt{W}$. Would you take this bet?
  \item \textbf{[5 Points]}  Can you propose a utility function such that you would take this bet? Verify that indeed taking the bet with your proposed utility function is better than not taking it.
\end{subquestions}
\question{Technology (10 Points)}
Calculate (i) the marginal products of both factors of production and (ii) the
technical rate of substitution between both factors for the following
production function:
\[
  f\left( x_1,x_2 \right)=x_1^{\frac{1}{4}}x_2^{\frac{1}{2}}
\]

\question{Cost Curves (5 Points)}
The marginal cost function of a firm is $MC\left( y \right)=2$. The firm has
fixed costs of 5. What is the total cost of producing $10$ units of output?


\question{Firm Supply and Industry Supply (20 Points)}
All firms operating in a perfectly competitive industry have the same cost function $c\left( y\right) = y^2 +1$. The market demand function is $D\left( p \right)=100-10p$.
\begin{subquestions}
  \item \textbf{[3 Points]} What is each firm's supply function?
  \item \textbf{[3 Points]} If there are 30 firms, what is the industry supply function?
  \item \textbf{[3 Points]} If there are 30 firms, what is the equilibrium price?
  \item \textbf{[3 Points]} If there are 30 firms, how much will each firm produce?
  \item \textbf{[3 Points]} If there are 30 firms, what is each firm's profits?
  \item \textbf{[5 Points]} How many firms will there be in the long run?
\end{subquestions}

\question{Equilibrium and Taxation (15 Points)}
The demand and supply functions for good are:
\begin{align*}
  D\left( p \right) &= 120 - p \\
  S\left( p \right) &= 2p
\end{align*}
\begin{subquestions}
  \item \textbf{[5 Points]} Calculate the equilibrium price and quantity.
\end{subquestions}
Suppose the government now adds a sales tax to the good of 20\%.
\begin{subquestions}
  \setcounter{enumi}{1}
  \item \textbf{[5 Points]} Calculate the price buyers pay, the price sellers receive and quantity sold after the tax is imposed.
  \item  \textbf{[5 Points]} What portion of the tax to buyers pay and what portion of the tax do
    sellers pay?
\end{subquestions}
\end{document}
